%% literature_review_ko.tex
\chapter{관련 연구}
\label{chap:litreview_ko}

\section{정보 검색의 기초}
\label{sec:lit_ir_ko}

정보 검색(Information Retrieval, IR)은 정보 요구에 응답하여 자료를 검색하는 시스템을 연구하는 학문 분야이다 \citep{manning2008ir}. 정형화된 검색 파이프라인은 문서 컬렉션을 색인화하고, 질의 시점에서 각 색인 문서를 질의에 대해 채점하여 순위 리스트를 생성하는 과정으로 구성된다. 현대 기업 검색 시스템은 수십 년간의 IR 연구, 특히 희소 단어 가중치와 밀집 표현 학습 분야의 성과를 기반으로 발전하였다.

\subsection{희소 검색: BM25}

Best Match 25(BM25) 검색 함수 \citep{robertson1995okapi,robertson2009bm25}는 실제 운영 시스템에서 가장 널리 사용되는 희소 검색 방법이다. BM25는 고전적인 TF-IDF 가중치 방식을 단어 빈도 포화 파라미터($k_1$)와 문서 길이 정규화 파라미터($b$)로 확장한다:

\begin{equation}
    \text{BM25}(d, q) = \sum_{t \in q} \text{IDF}(t) \cdot \frac{f(t,d) \cdot (k_1 + 1)}{f(t,d) + k_1 \cdot \left(1 - b + b \cdot \frac{|d|}{\text{avgdl}}\right)}
    \label{eq:bm25_ko}
\end{equation}

여기서 $f(t,d)$는 문서 $d$에서 단어 $t$의 빈도, $|d|$는 문서 길이, $\text{avgdl}$은 코퍼스의 평균 문서 길이이다. BM25는 역색인 조회를 통한 계산 효율성, 해석 가능성, 개체명·버전 번호·정책 코드 등 정확한 매칭 질의에 대한 강점을 지닌다. 주요 약점은 어휘 불일치(vocabulary mismatch)이다: 동일한 개념을 다른 어휘로 표현하는 질의와 문서는 낮은 BM25 점수를 받는다.

구조적(정책 코드, API 엔드포인트) 내용과 비구조적(개념 설명, 회의록) 내용이 혼재하는 기업 코퍼스에서 BM25의 강점과 약점은 모두 두드러지게 나타난다 \citep{robertson2004understanding}.

\subsection{밀집 검색}

밀집 검색(Dense Retrieval)은 질의와 문서를 연속 임베딩 공간의 밀집 벡터로 표현하고, 벡터 유사도로 관련성을 계산한다 \citep{karpukhin2020dpr}. 바이인코더 아키텍처를 사용하는 DPR은 개방형 질의응답에서 BM25를 크게 능가함을 보였으며, 특히 답변이 질의와 어휘적으로 다르게 표현된 경우에 효과적이다.

현대 밀집 검색 시스템은 검색 과제에 미세조정된 사전 훈련 트랜스포머 인코더를 사용한다. BAAI 일반 임베딩(BGE) 시리즈 \citep{xiao2023bge}는 다양한 도메인에서 강력한 검색 성능을 달성하는 경량 고품질 인코더를 제공한다. 본 연구에서 사용한 BGE-small-en-v1.5 모델은 텍스트를 384차원 벡터로 인코딩하며, 코사인 유사도 기반 검색을 지원한다.

그러나 밀집 검색에는 알려진 실패 유형이 존재한다. 일반 텍스트 코퍼스에서 훈련된 정적 밀집 임베딩은 날짜 순서나 버전 이력과 같은 시간적 순서 신호를 인코딩하지 못한다 \citep{manning2008ir}. 또한 임베딩 공간이 유사하지만 기능적으로 구별되는 개념들을 혼동할 수 있는 희귀 엔티티 처리에도 어려움을 보인다.

ColBERT \citep{khattab2020colbert}와 SPLADE \citep{formal2021splade}는 단일 모델 내에서 어휘적 매칭과 의미적 매칭의 상보적 강점을 결합하려는 희소-밀집 하이브리드 접근법이다. 그러나 이러한 방법들은 도메인 특화 훈련 데이터와 전문화된 색인 구조를 필요로 하므로, 임의의 기업 코퍼스에 즉시 적용하기 어렵다.

\subsection{하이브리드 검색과 상호 순위 융합}

희소 신호와 밀집 신호의 상보성을 고려할 때, 두 방법을 결합하는 하이브리드 검색 시스템이 기업 검색에서 널리 채택되었다 \citep{luan2021sparse_dense}. 상호 순위 융합(Reciprocal Rank Fusion, RRF) 알고리즘 \citep{cormack2009rrf}은 복수 검색 시스템의 순위 리스트를 결합하는 간단하고 파라미터가 적은 방법이다:

\begin{equation}
    \text{RRF}(d) = \sum_{r \in R} \frac{1}{k + r(d)}
    \label{eq:rrf_ko}
\end{equation}

여기서 $R$은 순위 리스트의 집합, $r(d)$는 리스트 $r$에서 문서 $d$의 순위, $k$는 상수이다(본 연구에서 $k=10$ 사용). RRF는 표준 IR 벤치마크에서 개별 순위 학습 방법과 콩도르세 투표 방식을 능가함이 입증되었다 \citep{cormack2009rrf}.

BEIR 벤치마크 \citep{thakur2021beir}는 생물의학, 법률, 일반 웹 검색 도메인을 포함한 다양한 검색 과제에서 하이브리드 검색이 강건함을 보인다는 증거를 제공한다.

\section{검색 증강 생성}
\label{sec:lit_rag_ko}

검색 증강 생성(RAG) \citep{lewis2020rag}은 검색된 구절에 조건화하여 언어 모델의 생성 출력을 생성함으로써, 모델이 매개변수에 인코딩되지 않은 도메인 특화 또는 최신 정보에 접근할 수 있도록 한다. RAG 아키텍처는 (1) 질의를 구절 집합으로 매핑하는 검색기와 (2) 질의 및 검색된 구절에 조건화하여 답변을 생성하는 생성기로 구성된다.

후속 연구들은 RAG를 여러 방향으로 확장하였다. Fusion-in-Decoder(FiD) \citep{izacard2021fid}는 각 검색 구절을 독립적으로 처리하여 디코더에서 표현을 융합함으로써 더 큰 검색 집합을 활용할 수 있다. RETRO \citep{borgeaud2022retro}는 검색을 사전 훈련 과정에 통합하여 수조 개의 토큰 규모 지식 검색을 가능하게 한다. Self-RAG \citep{asai2023selfrag}는 검색 시점과 검색된 구절 평가 방법을 결정하는 학습된 성찰 토큰을 생성기에 추가한다.

능동적 검색 증강 생성(Active RAG) \citep{jiang2023active}은 단일 검색 단계로 충분하지 않은 다단계 추론 시나리오를 다룬다. 이 접근법은 중간 생성 출력에 기반하여 새로운 검색 질의를 발행하는 생성-검색 교차 방식을 사용하며, 여러 문서에서 증거를 합성해야 하는 다중 홉 질의에 특히 유용하다.

\section{질의 분류와 적응형 검색}
\label{sec:lit_adaptive_ko}

질의 분류는 질의 특성에 기반하여 입력 질의를 사전 정의된 범주 또는 난이도 클래스 중 하나에 할당하는 과제이다. 고전적 IR에서 질의 분류는 특화된 검색 백엔드로 질의를 라우팅하거나 질의 확장 전략을 선택하는 데 적용되었다 \citep{baeza1999modern}. RAG 맥락에서 질의 분류는 주어진 질의에 대해 고품질 결과를 반환할 가능성이 높은 검색 전략을 선택하는 \emph{적응형 검색}을 가능하게 한다.

도메인 매칭 사전 훈련과 미세조정은 특화된 코퍼스에 대한 밀집 검색 성능을 크게 향상시킬 수 있음이 알려져 있다 \citep{oguz2021domain}. 그러나 미세조정은 비용이 많이 들고 레이블 데이터가 필요하다. 더 간단한 대안은 구조적 특성(질의 길이, 개체명 존재 여부, 추론 유형 등)에 기반하여 질의를 분류하고 결정론적 라우팅 규칙을 적용하는 것이다.

LLM 기반 질의 계획은 연쇄 추론 프롬프팅 \citep{wei2022chain}과 도구 사용 \citep{schick2023toolformer, yao2023react}의 성공을 기반으로, 언어 모델이 복잡한 질의를 분해하고, 검색 전략을 선택하며, 목표 지향적 검색 요청을 발행할 수 있도록 한다. WebGPT \citep{nakano2022webgpt}는 언어 모델이 웹 검색 도구를 효과적으로 사용하도록 훈련될 수 있음을 보였다. ReAct \citep{yao2023react}는 추론 트레이스와 액션 실행을 교차하는 방식으로 다단계 검색 증강 추론을 가능하게 한다.

\section{LLM 기반 계획 및 증강 언어 모델}
\label{sec:lit_llm_planning_ko}

증강 언어 모델 조사 연구 \citep{mialon2023augmented}는 검색, 도구 사용, 추론을 언어 모델 증강의 세 가지 주요 차원으로 식별한다. 검색 증강은 사실적 근거를 제공하고, 도구 사용은 모델의 행동 공간을 확장하며, 추론 증강(연쇄 추론, 스크래치패드, 자기 성찰)은 복합적 문제 해결을 향상시킨다.

대규모 언어 모델은 적절한 프롬프팅이 제공될 때 구조화된 분류 과제에서 강력한 제로샷 성능을 보인다 \citep{wei2022chain}. 구조화 프롬프트 접근법 \citep{mialon2023augmented}은 분류 과제에 대한 명시적 메타데이터를 프롬프트에 제공하여, 모델이 문맥 내 예시 없이도 의사결정 논리를 적용할 수 있도록 한다.

그러나 LLM은 정확한 논리 규칙이 필요한 과제에서 알려진 실패 유형을 보인다. Shi et al. \citep{shi2023large}은 LLM이 관련 없는 맥락에 쉽게 산만해지며, 혼동 정보가 포함된 경우 추론 과제 성능이 저하된다는 것을 보였다. 이는 검색 계획 맥락에서 중요한 의미를 가진다. 최적 전략이 특정 비명시적 특성(예: 정확한 식별자의 존재로 BM25 라우팅 필요)에 의존하는 경우, 모델의 학습된 사전 확률에 의해 이러한 규칙이 무시될 수 있다.

GPT-5는 가시적 출력 생성 전에 내부 연쇄 추론 계산(추론 토큰)을 사용하는 확장 추론 모델이다. 이 아키텍처는 \texttt{max\_completion\_tokens} 파라미터의 신중한 설정이 필요하다. 해당 파라미터는 추론 토큰과 가시적 출력 토큰 모두에 대한 공유 예산을 제공한다. 예산이 부족하면 모델이 추론 예산을 소진한 후에도 출력을 생성하지 못하는 문제가 발생한다. 이는 추론 모델 특유의 실패 유형으로, 본 연구에서 파악하고 해결하였다.

\section{정보 검색 시스템 평가}
\label{sec:lit_eval_ko}

표준 IR 평가 지표로는 Recall@$k$(상위 $k$개 결과에서 발견된 관련 문서의 비율), 평균 역순위(MRR; 첫 번째 관련 결과의 순위 역수의 평균), 정규화 할인 누적 이득(nDCG; 이상적 순위에 의해 정규화된 위치 할인 이득), Hit Rate@$k$(상위 $k$개에 최소 하나의 관련 문서가 포함되는지 여부) 등이 있다 \citep{manning2008ir, jaervelin2002ndcg}.

텍스트 검색 컨퍼런스(TREC) \citep{voorhees2005trec}는 1992년 이후 표준화된 IR 평가의 주요 동력이 되어 왔으며, 대규모 검색 벤치마크에서 사용되는 풀링 및 관련성 판단 방법론을 확립하였다. BEIR 벤치마크 \citep{thakur2021beir}는 18개 다양한 데이터셋에 걸쳐 밀집 및 하이브리드 검색 평가로 이 전통을 확장한다.

\section{소결}
\label{sec:lit_summary_ko}

관련 연구 검토를 통해 다음과 같은 결론을 도출할 수 있다: (1) BM25와 밀집 검색은 경쟁적이 아닌 상보적 방법이다; (2) Hybrid RRF는 이론적으로 충분히 뒷받침되고 광범위하게 검증된 신호 결합 접근법이다; (3) 적응형 검색 계획은 이론적으로 충분히 동기화되었으나 기업 코퍼스에 대한 실증적 검증이 부족하다; (4) LLM 기반 계획은 의미적으로 복잡한 질의에서 잠재적 강점을 제공하지만, 정확한 논리 규칙이 필요한 과제에서 알려진 실패 유형을 가진다.

본 논문은 기업 지식 관리의 적응형 검색 계획에 대한 이론적 동기와 실증적 검증 사이의 공백을 해소하기 위해, 목적에 맞게 구축된 기업 지식 코퍼스에서 5가지 접근법(BM25, Dense, Hybrid, 규칙 계획기, LLM 계획기) 모두를 최초로 통제된 방식으로 비교·평가한다.
